\begin{figure*}[h]
  \centering
  {\Large\bfseries\sffamily Resubmission of SIGMETRICS-26-WINTER Submission \#210}
  % \vspace{2pt}
  % \caption*{}  
\end{figure*}

We sincerely thank all reviewers and the PC for their comprehensive and insightful comments.
We have carefully revised the paper following their constructive suggestions.
In this appendix, Section~\ref{sec:major-revisions} summarizes the revisions made in response to the major comments, Section~\ref{sec:specific-replies} provides specific replies to each reviewer, and Section~\ref{expr:ci_tables} reports supplementary confidence-interval tables for the revised experiments.

\section{Revisions Following Major Comments}
\label{sec:major-revisions}

\subsection{Clarify the contributions, especially on the implementation side}
\label{subsec:rev-contributions}
We added implementation details for X-Blossom-Pro and X-Blossom++ to the first two items in the contribution list in Subsection~\ref{subsec:intro-exec-summary}.

\subsection{Include more details of the experimental setup}
\label{subsec:rev-exp-setup}
We expanded the experimental setup description to include the cloud platform, operating system, instance types, processor and accelerator models, CPU core/thread counts, and the framework/threading model used by each implementation. 
Specifically, we clarified that XB-Pro uses C++ \texttt{std::thread} with pthread support and XB++ uses CUDA kernels; for the Ligra and Gunrock baselines, we state that Ligra uses OpenMP for CPU parallelism and Gunrock runs on a CUDA-based graph analytics framework.
These additions provide a more complete and reproducible description of the hardware and software environment used in the revised experiments. Further details can be found in Subsubsection~\ref{subsubsec:eval-exp-env} and Subsection~\ref{subsec:baselines}.

\subsection{Scalability test}
\label{subsec:rev-scalability}
We added a new scalability experiment that evaluates both XB-Pro and XB++ under node-level and edge-level load-balancing strategies on three graph datasets. 
The results show that edge-level load balancing generally achieves higher speedup than node-level load balancing, and that GPU scaling is higher than CPU scaling due to lower CUDA-thread overhead and efficient SM-level parallel execution.
Further details can be found in the Subsection~\ref{subsec:scalability_test}.

\subsection{Clarify the i9 cores vs. threads}
\label{subsec:rev-cores-threads}
We revised the hardware description to explicitly distinguish physical CPU cores from hardware threads. In the updated setup, the CPU instance is reported as having 48 physical cores and 96 hardware threads, avoiding ambiguity between core count and thread count. Further details can be found in Subsubsection~\ref{subsubsec:eval-exp-env}.

\subsection{Add confidence intervals}
\label{subsec:rev-confidence-intervals}
We added confidence intervals (CIs) to Tables~\ref{tab:graph-metrics} and \ref{tab:reuse} to demonstrate the stability of the reported averages.
For Table~\ref{tab:pe}, Table~\ref{tab:memory}, and Figures~\ref{plot:load-balance}, \ref{plot:scalability_test}, \ref{plot:runtimes}, and \ref{plot:inst_ratio}, we provide the corresponding CI tables in the appendix (Section~\ref{expr:ci_tables}) to keep the main presentation compact and readable.

\subsection{Include details about execution phases}
\label{subsec:rev-execution-phases}
We revised the last two paragraphs of the load-balancing discussion in Subsubsection~\ref{subsubsec:exp_load_balancing}; the revised sequential-work paragraph clarifies which parts of X-Blossom++ remain sequential or expensive~\hyperref[para:xbpp-sequential-work]{(see the sequential-work paragraph)}, and the revised parallel-work paragraph clarifies which major search phases are parallel~\hyperref[para:xbpp-parallel-work]{(see the parallel-work paragraph)}. Together, the revised text explains that CPU-to-GPU speedup is governed by the fraction of work that can be exposed as fine-grained parallel GPU tasks versus the remaining sequential blossom/path-table work, so it does not scale directly with the raw increase in GPU thread count.

\subsection{Run CPU and GPU experiments on comparable platforms}
\label{subsec:rev-fair-platform}
To address this concern, we moved the experiments from the previous local i9/A100 setup to server-class Amazon Web Services (AWS) instances. Although we still used two separate instances because AWS GPU instances do not expose the bare-metal memory-system information required for our CPU baseline evaluation, the comparison is fair because both platforms are server-class AWS instances selected for the corresponding CPU and GPU baselines rather than consumer or mobile hardware. The CPU baselines were evaluated on an m7i.metal-24xl bare-metal instance, a current-generation AWS server-class instance with 4th Generation Intel Xeon processors, which provides direct access to the memory-system information required by our CPU-side analysis. The GPU baselines were evaluated on a server-class AWS GPU instance using an NVIDIA RTX PRO 6000 Blackwell Server Edition GPU, which is newer than the A100-class platform and provides competitive or better performance for the evaluated GPU baselines. Further details can be found in Subsubsection~\ref{subsubsec:eval-exp-env}.

\subsection{Explain why the ratio is small for some larger graphs in Table~5}
\label{subsec:rev-inst-ratio}
This issue no longer appears after moving the experiments from the previous A100 platform to the RTX PRO 6000 Blackwell platform. The A100 exposes 108 SMs, whereas the RTX PRO 6000 Blackwell exposes 188 SMs, providing substantially more SM-level execution resources. With the revised platform, XB++ achieves a higher instruction issue rate than BFS-Gunrock on all evaluated graphs, which is consistent with our conclusion that XB++ exposes more GPU parallelism through fine-grained edge-level work. We updated the first sentence in 
\hyperref[para:inst_gpu]{this paragraph} to explicitly state this observation and connect it to our parallelism-exploitation conclusion.

\subsection{Comment on the variability in the improvement of the speedup across different graphs in Table~\ref{tab:reuse}}
\label{subsec:rev-reuse-variability}
We added one dynamic metric, ReuseRatio, to Table~\ref{tab:reuse} and used it to explain the variability in reuse speedups across graphs in Subsection~\ref{subsub:exp_reuse}. 
ReuseRatio measures the percentage of alternating-tree nodes preserved by the reuse mechanism rather than reset, so it directly captures the internal behavior that alternating-tree reuse is designed to exploit. 
The revised discussion shows that datasets with higher ReuseRatio generally obtain larger speedups.

\subsection{Expand a bit on how SIMT execution helps to keep up massive concurrency}
\label{subsec:rev-simt-concurrency}
We expanded the explanation of SIMT execution to clarify how it helps XB++ maintain massive concurrency and reduce effective memory and control-flow latency, especially when edge-level load balancing exposes enough fine-grained tasks in \hyperref[para:smit_execution]{this paragraph}.

\subsection{Explain small working set a bit more on Page~17}
\label{subsec:rev-working-set}
We added a definition of working set to the data-movement discussion in Subsection~\ref{subsec:profiling_method}.
